\section{Copper-loss ellipsoid and minimum-loss operation}

The copper-loss model is
\begin{equation}
 \Pcu=\frac32\Rs(\id^2+\iq^2)+\Rexc\iexc^2
 =\state\trans\Rl\state,
 \qquad
 \Rl=\diag\left(\frac32\Rs,\frac32\Rs,\Rexc\right).
 \label{eq:copper-loss}
\end{equation}
For $R_s>0$ and $R_e>0$, $\Rl$ is positive definite.  Hence
$\state\trans\Rl\state=P_{\mathrm{Cu,max}}$ is a closed ellipsoid, with
semi-axes
\begin{equation}
 \sqrt{\frac{\Pmax}{(3/2)\Rs}},\quad
 \sqrt{\frac{\Pmax}{(3/2)\Rs}},\quad
 \sqrt{\frac{\Pmax}{\Rexc}}.
\end{equation}
This is a copper-loss or thermal-budget ellipsoid, not an output-power
ellipsoid.  Its definiteness sharply contrasts with the rank-deficient voltage
cylinder.

A physically meaningful EESM analogue of MTPA is
\begin{equation}
 \min_{\state}\Pcu\quad\text{subject to}\quad\Te=T^*.
 \label{eq:min-loss}
\end{equation}
Why should gradients appear?  Imagine expanding a copper-loss ellipsoid from
the origin until it first reaches the requested torque surface.  At first
contact the two level surfaces share a tangent plane.  Their normal vectors
must therefore be parallel, and gradients are precisely those normals.  This
geometric contact condition gives
$\nabla\Pcu=\mu\nabla\Te$ at a regular solution, i.e.,
\begin{align}
 3\Rs\id&=\mu\kt\dL\iq,\\
 3\Rs\iq&=\mu\kt(\Lexc\iexc+\dL\id),\\
 2\Rexc\iexc&=\mu\kt\Lexc\iq.
\end{align}
On a nonzero-torque branch, the first and third equations give
\begin{equation}
 \boxed{\id=\frac{2\Rexc\dL}{3\Rs\Lexc}\iexc.}\label{eq:minloss-ratio}
\end{equation}
The remaining equations determine a constant ratio for $i_q/i_e$.  This ray
property follows from homogeneity:
\begin{equation}
 \Te(\rho\state)=\rho^2\Te(\state),\qquad
 \Pcu(\rho\state)=\rho^2\Pcu(\state).
\end{equation}
The optimum direction is independent of requested torque; only scale changes.

Equivalently, maximizing torque for a fixed loss budget gives
\begin{equation}
 \Hm\state=\lambda\Rl\state,\label{eq:generalized-eigenproblem}
\end{equation}
but jumping directly to the phrase generalized eigenvalue problem hides its
geometry.  The loss constraint
$\state\trans\Rl\state=P_0$ is an ellipsoid because equal Euclidean current
steps have unequal copper cost.  Define the loss-whitened coordinate
\begin{equation}
 \vect y=\Rl^{1/2}\state,\qquad
 \state=\Rl^{-1/2}\vect y.
\end{equation}
Then explicitly
\begin{equation}
 \state\trans\Rl\state
 =\vect y\trans\Rl^{-1/2}\Rl\Rl^{-1/2}\vect y
 =\vect y\trans\vect y=P_0.
\end{equation}
Whitening therefore changes coordinates so that one unit of Euclidean
distance represents one unit of copper-loss resource in every direction.  It
exposes fair resource comparison, while hiding the original amperes until the
result is transformed back.  Torque becomes
\begin{equation}
 \Te=\vect y\trans(\Rl^{-1/2}\Hm\Rl^{-1/2})\vect y.
\end{equation}
On the sphere, the best direction is the largest-eigenvalue direction of the
symmetric whitened torque matrix:
\begin{equation}
 (\Rl^{-1/2}\Hm\Rl^{-1/2})\vect y=\lambda\vect y.
\end{equation}
Substituting $\vect y=\Rl^{1/2}\state$ and multiplying by
$\Rl^{1/2}$ recovers
$\Hm\state=\lambda\Rl\state$.  The generalized eigenproblem is thus an
ordinary eigenproblem viewed back in physical current coordinates.

\paragraph{Synthesis.}
The loss ellipsoid supplies a resource metric, tangency supplies the optimum
condition, whitening makes that metric Euclidean, and the dominant eigenvector
selects the most torque-effective direction per unit loss.  This argument
requires $\Rl\succ0$; if a resistance vanished, whitening would be singular
and the corresponding current direction would have zero modeled loss.
